% Original teaching diagram from the exact example in Chapter 19.
\begin{tikzpicture}[x=1mm,y=1mm,>=Stealth,font=\figurefont\fontsize{8}{10}\selectfont,
 v/.style={circle,draw=FigureStroke,fill=white,minimum size=4.8mm,inner sep=0pt},
 e/.style={draw=FigureStroke,line width=.7pt},
 op/.style={rounded corners=1.8pt,minimum width=43mm,minimum height=8mm,align=center},
 sub/.style={font=\FigureNoteFont\fontsize{7.5}{9}\selectfont}]
  \foreach \shift/\name in {0/\mathcal G,54/\mathcal H} {
    \begin{scope}[xshift=\shift mm]
      \coordinate (p) at (2,38);
      \coordinate (q) at (20,38);
      \coordinate (r) at (20,23);
      \coordinate (s) at (2,23);
      \draw[e] (p)--(q)--(r)--(s)--cycle;
      \ifnum\shift=0 \draw[FigureBlue,line width=1.1pt] (p)--(r);\fi
      \foreach \vertex/\idx in {p/1,q/2,r/3,s/4} \node[v] at (\vertex) {$\idx$};
      \node[above=0.5mm] at (11,41) {$\name$};
    \end{scope}
  }
  \draw[->,FigureMuted] (26,30.5)--node[above=1pt] {删去 $13$} (50,30.5);
  \node[sub] at (38,15) {(a)};
  \node[op,draw=FigureGreen,fill=FigureGreenFill] (sparse) at (11,2)
    {$\bm M\bm z=\bm r^{(0)}$\\求方向 $\bm z$};
  \node[op,draw=FigureBlue,fill=FigureBlueFill] (orig) at (65,2)
    {$\bm z^\top\bm L_G\bm z$\\用于确定步长};
  \draw[->,FigureStroke] (sparse)--(orig);
  \node[sub] at (38,-7) {(b)};
\end{tikzpicture}
